\begin{proposition}[2.5.5]
\label{IV.2.5.5-fr}
Sean $f:X\to Y$ un morfismo, $\mathscr F$ un $\mathcal O_X$-módulo
cuasicoherente de tipo finito y plano respecto de $f$, y $\mathscr E$ un
$\mathcal O_Y$-módulo cuasicoherente de tipo finito; para todo $y\in Y$,
se pone $\mathscr F_y=\mathscr F\otimes_{\mathcal O_Y}\boldsymbol{k}(y)$. Para que un punto
$x\in X$ sea punto maximal de
$\operatorname{Supp}(\mathscr E\otimes_Y\mathscr F)$, es necesario y suficiente
que $y=f(x)$ sea punto maximal de $\operatorname{Supp}(\mathscr E)$ y que $x$
sea punto maximal de $\operatorname{Supp}(\mathscr F_y)$ en $f^{-1}(y)$. Si
así ocurre, se tiene
\end{proposition}

\begin{env}[2.5.5.1]
\label{IV.2.5.5.1-fr}
\[
\operatorname{long}\bigl((\mathscr E\otimes_Y\mathscr F)_x\bigr)
=\operatorname{long}(\mathscr E_y).
 \operatorname{long}\bigl((\mathscr F_y)_x\bigr)
\]
\end{env}

Es claro que
$f(\operatorname{Supp}(\mathscr E\otimes_Y\mathscr F))
\subset\operatorname{Supp}(\mathscr E)$ $(\mathbf{0}_{\mathrm I},5.2.2)$; la imagen por
$f$ de toda componente irreducible de
$\operatorname{Supp}(\mathscr E\otimes_Y\mathscr F)$ está, por tanto, contenida
en una componente irreducible de $\operatorname{Supp}(\mathscr E)$. Es posible
limitarse al caso en que $\operatorname{Supp}(\mathscr F)=X$. En efecto
(\textbf{Err}\textsubscript{III}, 30), puesto que $\mathscr F$ es de
tipo finito, existen un subpreesquema cerrado $X'$ de $X$ cuyo espacio
subyacente es $\operatorname{Supp}(\mathscr F)$ y un $\mathcal O_{X'}$-módulo
cuasicoherente de tipo finito $\mathscr G$ tales que, si $j:X'\to X$ es la
inyección canónica, se tiene $\mathscr F=j_*(\mathscr G)$. Si se pone
$f'=f\circ j$, es claro que $\mathscr G$ es plano respecto de $f'$ y que se
tiene
$\operatorname{Supp}(\mathscr E\otimes_Y\mathscr F)
=\operatorname{Supp}(\mathscr E\otimes_Y\mathscr G)$.
Supongamos, pues, que $\operatorname{Supp}(\mathscr F)=X$; si $Z$ es una
componente irreducible de $\operatorname{Supp}(\mathscr E)$, se tiene
$f^{-1}(Z)\subset\operatorname{Supp}(\mathscr E\otimes_Y\mathscr F)$
(\textbf{I}, 9.1.13),
y de (2.3.4) se sigue que toda componente irreducible de $f^{-1}(Z)$ domina
$Z$; en otras palabras, si $x$ es el punto genérico de una componente
irreducible de $\operatorname{Supp}(\mathscr E\otimes_Y\mathscr F)$ contenida
en $f^{-1}(Z)$, $y=f(x)$ es el punto genérico de $Z$; además,
$(\mathbf{0}_{\mathrm I},2.1.8)$, $x$ es el punto genérico de una componente irreducible
de $f^{-1}(y)=\operatorname{Supp}(\mathscr F_y)$ (\textbf I, 9.1.13), y recíprocamente
todo punto genérico de una de estas componentes es punto genérico de una
componente irreducible de $f^{-1}(Z)$.

\oldpage[IV-2]{24}
Queda por demostrar (2.5.5.1); se tiene
$(\mathscr E\otimes_Y\mathscr F)_x
=\mathscr E_y\otimes_{\mathcal O_y}\mathscr F_x$
(\textbf{I}, 9.1.12)
y
$(\mathscr F_y)_x=\mathscr F_x/\mathfrak m_y\mathscr F_x$; la cuestión se reduce, por
tanto, a demostrar el

\begin{lemma}[2.5.5.2]
\label{IV.2.5.5.2-fr}
Sean $A$ y $B$ dos anillos locales, $\rho:A\to B$ un homomorfismo local y
$\mathfrak m$ el ideal maximal de $A$. Sean $M$ un $A$-módulo y $N$ un
$B$-módulo que es un $A$-módulo fielmente plano y tal que
$N/\mathfrak mN$ es un $B$-módulo de longitud finita; entonces se tiene
\end{lemma}

\begin{env}[2.5.5.3]
\label{IV.2.5.5.3-fr}
\[
\operatorname{long}_B(M\otimes_A N)
=\operatorname{long}_A(M).
 \operatorname{long}_B(N/\mathfrak mN).
\]
\end{env}

Si $M$ es de longitud infinita, lo mismo ocurre con $M\otimes_A N$, pues para
toda sucesión estrictamente creciente de $n$ submódulos $M_i$ de $M$
($1\leq i\leq n$), los $M_i\otimes_A N$ se identifican con sub-$B$-módulos de
$M\otimes_A N$ distintos dos a dos; como $N\ne\mathfrak mN$ puesto que $N$ es
un $A$-módulo fielmente plano, la fórmula (2.5.5.3) es cierta en este caso.
Supongamos, pues, que $M$ es de longitud finita. Si $M=0$, los dos miembros de
(2.5.5.3) son nulos, por lo que se puede suponer $M\ne0$. Los
$\mathfrak m^kM$ son submódulos de $M$, y por tanto de longitud finita, y si
$\mathfrak m^kM\ne0$, del lema de Nakayama se sigue que
$\mathfrak m^{k+1}M\ne\mathfrak m^kM$; existe, pues, necesariamente un entero
$r$ tal que $\mathfrak m^rM=0$. Los $\mathfrak m^kM\otimes_A N$ se identifican
entonces con sub-$B$-módulos de $M\otimes_A N$,
$(\mathfrak m^kM\otimes_A N)/(\mathfrak m^{k+1}M\otimes_A N)$
siendo isomorfo a
$(\mathfrak m^kM/\mathfrak m^{k+1}M)\otimes_A N$,
y por tanto también a
$(\mathfrak m^kM/\mathfrak m^{k+1}M)
\otimes_{A/\mathfrak m}(N/\mathfrak mN)$
$(0\leq k\leq r\setmin 1)$~;
la longitud de este último, considerado como $B$-módulo, es, por tanto, el
producto de $\operatorname{long}_B(N/\mathfrak mN)$ por el rango del espacio
vectorial sobre $(A/\mathfrak m)$
$\mathfrak m^kM/\mathfrak m^{k+1}M$, rango que es también la longitud del
$A$-módulo $\mathfrak m^kM/\mathfrak m^{k+1}M$. Sumando (para
$0\leq k\leq r\setmin 1$), se deduce inmediatamente la fórmula (2.5.5.3).

\begin{remark}[2.5.5.4]
\label{IV.2.5.5.4-fr}
Se observará que, cuando $N$ es un $A$-módulo \emph{de tipo finito}, decir que
$N$ es un $A$-módulo fielmente plano equivale a decir que $N\ne0$ y que $N$
es un $A$-módulo \emph{plano}; en efecto, el lema de Nakayama muestra entonces que
$\mathfrak mN\ne N$.
\end{remark}
